\section{自由李代数}

我们还须谈谈一系列关于李代数的工作，它们与李群理论的联系十分微弱；但另一方面，这些研究在“抽象”群理论中、尤其是在幂零群理论中有着重要的应用。

这件事的起源是 P. 霍尔发表于 1932 年的工作{[}144{]}。不过那里完全谈不上李代数：P. 霍尔着眼于研究某一类 p 群，即他称为“正则”的那类 p 群。但这导致他仔细考察了一个群的迭代换位子和下中心列；他借此机会建立了雅可比恒等式的一个变体，以及“霍尔公式”

\[
(x y) ^ {n} = x ^ {n} y ^ {n} (x, y) ^ {n (n - 1) / 2} \dots
\]

不久之后（1935—1937 年），W. 马格努斯{[}215 a 和 b{]}与 E. 维特{[}337 b{]}的基本工作问世。在{[}215 a{]}中，马格努斯使用了与豪斯多夫相同的形式级数代数 \(\hat{A}\)（后来称为“马格努斯代数”）；他把自由群 \(F\) 嵌入其中，并利用 \(\hat{A}\) 的自然过滤，得到了一个递降序列 \((F_n)\)，其中的群都是 \(F\) 的子群；这是最早的过滤例子之一。他猜想 \(F_n\) 与 \(F\) 的下中心列的各项相重合。这一猜想在他的第二篇论文{[}215 b{]}中得到证明；也正是在那里，他明确地把他的思想与 P. 霍尔的思想联系起来，并且定义了自由李代数 \(L\)（作为 \(A\) 的一个子代数），而且实质上证明了它与分次的 \(F\) 相等同。在{[}337 b{]}中，维特在若干方面完善了这一结果。他特别证明了 \(L\) 的包络代数是一个自由结合代数，并由此立即导出 \(L\) 的齐次分量的秩（“维特公式”）。

至于以“霍尔基”之名著称的 L 的基的确定，看来它直到 1950 年才第一次出现，即出现在 M. 霍尔的一篇短文{[}143{]}中，尽管它已隐含在上文引述的 P. 霍尔与 W. 马格努斯的工作里。

